% Options for packages loaded elsewhere
\PassOptionsToPackage{unicode}{hyperref}
\PassOptionsToPackage{hyphens}{url}
\documentclass[
  10pt,
]{article}
\usepackage{xcolor}
\usepackage[margin=1.6cm]{geometry}
\usepackage{amsmath,amssymb}
\setcounter{secnumdepth}{-\maxdimen} % remove section numbering
\usepackage{iftex}
\ifPDFTeX
  \usepackage[T1]{fontenc}
  \usepackage[utf8]{inputenc}
  \usepackage{textcomp} % provide euro and other symbols
\else % if luatex or xetex
  \usepackage{unicode-math} % this also loads fontspec
  \defaultfontfeatures{Scale=MatchLowercase}
  \defaultfontfeatures[\rmfamily]{Ligatures=TeX,Scale=1}
\fi
\usepackage{lmodern}
\ifPDFTeX\else
  % xetex/luatex font selection
  \setmainfont[]{Arial Unicode MS}
  \setmonofont[]{Menlo}
\fi
% Use upquote if available, for straight quotes in verbatim environments
\IfFileExists{upquote.sty}{\usepackage{upquote}}{}
\IfFileExists{microtype.sty}{% use microtype if available
  \usepackage[]{microtype}
  \UseMicrotypeSet[protrusion]{basicmath} % disable protrusion for tt fonts
}{}
\makeatletter
\@ifundefined{KOMAClassName}{% if non-KOMA class
  \IfFileExists{parskip.sty}{%
    \usepackage{parskip}
  }{% else
    \setlength{\parindent}{0pt}
    \setlength{\parskip}{6pt plus 2pt minus 1pt}}
}{% if KOMA class
  \KOMAoptions{parskip=half}}
\makeatother
\usepackage{longtable,booktabs,array}
\newcounter{none} % for unnumbered tables
\usepackage{calc} % for calculating minipage widths
% Correct order of tables after \paragraph or \subparagraph
\usepackage{etoolbox}
\makeatletter
\patchcmd\longtable{\par}{\if@noskipsec\mbox{}\fi\par}{}{}
\makeatother
% Allow footnotes in longtable head/foot
\IfFileExists{footnotehyper.sty}{\usepackage{footnotehyper}}{\usepackage{footnote}}
\makesavenoteenv{longtable}
\setlength{\emergencystretch}{3em} % prevent overfull lines
\providecommand{\tightlist}{%
  \setlength{\itemsep}{0pt}\setlength{\parskip}{0pt}}
\usepackage{xurl}
\urlstyle{same}
\setlength{\emergencystretch}{3em}
\sloppy
\usepackage{bookmark}
\IfFileExists{xurl.sty}{\usepackage{xurl}}{} % add URL line breaks if available
\urlstyle{same}
\hypersetup{
  hidelinks,
  pdfcreator={LaTeX via pandoc}}

\author{}
\date{}

\begin{document}

\section{JMH System Report}\label{jmh-system-report}

\begin{itemize}
\tightlist
\item
  Generated: \texttt{2026-03-15\ 01:19:08}
\item
  Command:
  \texttt{hw01/tools/jmh\_run\_with\_system\_metrics.py\ -\/-project-root\ .\ -\/-out-dir\ hw01/target/jmh-system-report\ -\/-rebuild-only}
\end{itemize}

\subsection{Summary}\label{summary}

\begin{itemize}
\tightlist
\item
  Benchmarks executed: \textbf{16}
\item
  Summary CSV: \texttt{operation\_summary.csv}
\item
  Flattened JMH CSV: \texttt{jmh\_results\_flat.csv}
\end{itemize}

\subsection{ExtendableHashTableBenchmark}\label{extendablehashtablebenchmark}

{\def\LTcaptype{none} % do not increment counter
\begin{longtable}[]{@{}
  >{\raggedright\arraybackslash}p{(\linewidth - 10\tabcolsep) * \real{0.1304}}
  >{\raggedleft\arraybackslash}p{(\linewidth - 10\tabcolsep) * \real{0.1739}}
  >{\raggedleft\arraybackslash}p{(\linewidth - 10\tabcolsep) * \real{0.1739}}
  >{\raggedleft\arraybackslash}p{(\linewidth - 10\tabcolsep) * \real{0.1739}}
  >{\raggedleft\arraybackslash}p{(\linewidth - 10\tabcolsep) * \real{0.1739}}
  >{\raggedleft\arraybackslash}p{(\linewidth - 10\tabcolsep) * \real{0.1739}}@{}}
\toprule\noalign{}
\begin{minipage}[b]{\linewidth}\raggedright
operation
\end{minipage} & \begin{minipage}[b]{\linewidth}\raggedleft
kind
\end{minipage} & \begin{minipage}[b]{\linewidth}\raggedleft
duration\_s
\end{minipage} & \begin{minipage}[b]{\linewidth}\raggedleft
p95\_cpu\_\%
\end{minipage} & \begin{minipage}[b]{\linewidth}\raggedleft
peak\_rss\_mb
\end{minipage} & \begin{minipage}[b]{\linewidth}\raggedleft
score\_median
\end{minipage} \\
\midrule\noalign{}
\endhead
\bottomrule\noalign{}
\endlastfoot
buildFromScratch & build & 106.776 & 0.800 & 334.703 & 413.194 \\
getHit & read/query & 653.918 & 0.200 & 331.516 & 26.637 \\
getMiss & read/query & 654.388 & 0.200 & 189.484 & 35.872 \\
putThenRemove & write & 654.354 & 0.200 & 186.562 & 16.956 \\
removeThenPut & write & 653.696 & 0.200 & 184.953 & 13.758 \\
updateExisting & write & 653.978 & 0.200 & 324.344 & 16.016 \\
\end{longtable}
}

\subsubsection{buildFromScratch}\label{buildfromscratch}

\begin{itemize}
\tightlist
\item
  Operation type: \texttt{build}
\item
  Parameters:

  \begin{itemize}
  \tightlist
  \item
    \texttt{bucketCapacity}: 8, 16
  \item
    \texttt{n}: 200000, 350000
  \item
    \texttt{seed}: 42
  \end{itemize}
\end{itemize}

\paragraph{Notes}\label{notes}

\begin{itemize}
\tightlist
\item
  Построение самое тяжелое в этом алгоритме, median \textasciitilde413
  ms/op, сильный разброс по параметрам \texttt{n},
  \texttt{bucketCapacity}.
\item
  Фикс: заранее резервировать структуру бакетов/директории и
  минимизировать перевыделения при split, переиспользовать рабочие
  буферы чтения/записи.
\end{itemize}

\subsubsection{getHit}\label{gethit}

\begin{itemize}
\tightlist
\item
  Operation type: \texttt{read/query}
\item
  Parameters:

  \begin{itemize}
  \tightlist
  \item
    \texttt{bucketCapacity}: 8, 16
  \item
    \texttt{n}: 200000, 350000
  \item
    \texttt{seed}: 42
  \end{itemize}
\end{itemize}

\paragraph{Notes}\label{notes-1}

\begin{itemize}
\tightlist
\item
  Чтение hit стабильно быстрое.
\item
  Фикс: кэш горячих бакетов в state бенчмарка для лучшей locality,
  убрать лишние чтения метаданных бакета в hot-path get.
\end{itemize}

\subsubsection{getMiss}\label{getmiss}

\begin{itemize}
\tightlist
\item
  Operation type: \texttt{read/query}
\item
  Parameters:

  \begin{itemize}
  \tightlist
  \item
    \texttt{bucketCapacity}: 8, 16
  \item
    \texttt{n}: 200000, 350000
  \item
    \texttt{seed}: 42
  \end{itemize}
\end{itemize}

\paragraph{Notes}\label{notes-2}

\begin{itemize}
\tightlist
\item
  Miss быстрее hit, что логично, так как раньше завершается без чтения
  значения.
\end{itemize}

\subsubsection{putThenRemove}\label{putthenremove}

\begin{itemize}
\tightlist
\item
  Operation type: \texttt{write}
\item
  Parameters:

  \begin{itemize}
  \tightlist
  \item
    \texttt{bucketCapacity}: 8, 16
  \item
    \texttt{n}: 100000, 200000
  \item
    \texttt{seed}: 42
  \end{itemize}
\end{itemize}

\paragraph{Notes}\label{notes-3}

\begin{itemize}
\tightlist
\item
  Запись+удаление: avgt \textasciitilde0.0260 us/op, thrpt
  \textasciitilde38.7 ops/us.
\item
  Фикс: объединить серию мелких записей в один проход, переиспользовать
  tombstone-слоты перед split бакета.
\end{itemize}

\subsubsection{removeThenPut}\label{removethenput}

\begin{itemize}
\tightlist
\item
  Operation type: \texttt{write}
\item
  Parameters:

  \begin{itemize}
  \tightlist
  \item
    \texttt{bucketCapacity}: 8, 16
  \item
    \texttt{n}: 100000, 200000
  \item
    \texttt{seed}: 42
  \end{itemize}
\end{itemize}

\paragraph{Notes}\label{notes-4}

\begin{itemize}
\tightlist
\item
  Сценарий remove-\textgreater put медленнее put-\textgreater remove.
\item
  Фикс: после remove сразу оставлять slot в состоянии, удобном для
  последующего put без повторного скана.
\end{itemize}

\subsubsection{updateExisting}\label{updateexisting}

\begin{itemize}
\tightlist
\item
  Operation type: \texttt{write}
\item
  Parameters:

  \begin{itemize}
  \tightlist
  \item
    \texttt{bucketCapacity}: 8, 16
  \item
    \texttt{n}: 100000, 200000
  \item
    \texttt{seed}: 42
  \end{itemize}
\end{itemize}

\paragraph{Notes}\label{notes-5}

\begin{itemize}
\tightlist
\item
  Update близок к другим write-операциям.
\item
  Фикс: если новое значение равно старому, выходить без физической
  записи.
\end{itemize}

\subsection{LshBenchmark}\label{lshbenchmark}

{\def\LTcaptype{none} % do not increment counter
\begin{longtable}[]{@{}
  >{\raggedright\arraybackslash}p{(\linewidth - 10\tabcolsep) * \real{0.1304}}
  >{\raggedleft\arraybackslash}p{(\linewidth - 10\tabcolsep) * \real{0.1739}}
  >{\raggedleft\arraybackslash}p{(\linewidth - 10\tabcolsep) * \real{0.1739}}
  >{\raggedleft\arraybackslash}p{(\linewidth - 10\tabcolsep) * \real{0.1739}}
  >{\raggedleft\arraybackslash}p{(\linewidth - 10\tabcolsep) * \real{0.1739}}
  >{\raggedleft\arraybackslash}p{(\linewidth - 10\tabcolsep) * \real{0.1739}}@{}}
\toprule\noalign{}
\begin{minipage}[b]{\linewidth}\raggedright
operation
\end{minipage} & \begin{minipage}[b]{\linewidth}\raggedleft
kind
\end{minipage} & \begin{minipage}[b]{\linewidth}\raggedleft
duration\_s
\end{minipage} & \begin{minipage}[b]{\linewidth}\raggedleft
p95\_cpu\_\%
\end{minipage} & \begin{minipage}[b]{\linewidth}\raggedleft
peak\_rss\_mb
\end{minipage} & \begin{minipage}[b]{\linewidth}\raggedleft
score\_median
\end{minipage} \\
\midrule\noalign{}
\endhead
\bottomrule\noalign{}
\endlastfoot
buildIndex & build & 488.716 & 0.200 & 2304.516 & 307.792 \\
candidatesQuery & read/query & 332.474 & 0.200 & 723.000 & 13.563 \\
nearDuplicatesFullScan & read/query & 483.371 & 0.200 & 324.719 &
425.400 \\
nearDuplicatesLsh & read/query & 371.598 & 0.200 & 1502.062 & 9.224 \\
addDocument & write & 214.170 & 0.210 & 2702.844 & 14.595 \\
\end{longtable}
}

\subsubsection{addDocument}\label{adddocument}

\begin{itemize}
\tightlist
\item
  Operation type: \texttt{write}
\item
  Parameters:

  \begin{itemize}
  \tightlist
  \item
    \texttt{baseDocs}: 20000
  \item
    \texttt{seed}: 42
  \item
    \texttt{wordsPerDoc}: 16
  \end{itemize}
\end{itemize}

\paragraph{Notes}\label{notes-6}

\begin{itemize}
\tightlist
\item
  Добавление документа быстрое, но общиее использование памяти большое
  на больших данных.
\item
  Фикс: переиспользовать временные массивы/хеши для шинглов между
  вызовами.
\end{itemize}

\subsubsection{buildIndex}\label{buildindex}

\begin{itemize}
\tightlist
\item
  Operation type: \texttt{build}
\item
  Parameters:

  \begin{itemize}
  \tightlist
  \item
    \texttt{docs}: 20000, 60000
  \item
    \texttt{seed}: 42
  \item
    \texttt{wordsPerDoc}: 16
  \end{itemize}
\end{itemize}

\paragraph{Notes}\label{notes-7}

\begin{itemize}
\tightlist
\item
  Самая тяжелая стадия LSH: avgt median \textasciitilde1250 ms/op, пик
  памяти \textasciitilde2.3 GB.
\item
  Фикс: перейти на примитивные коллекции/массивы для сигнатур.
\end{itemize}

\subsubsection{candidatesQuery}\label{candidatesquery}

\begin{itemize}
\tightlist
\item
  Operation type: \texttt{read/query}
\item
  Parameters:

  \begin{itemize}
  \tightlist
  \item
    \texttt{docs}: 20000, 60000
  \item
    \texttt{seed}: 42
  \item
    \texttt{wordsPerDoc}: 16
  \end{itemize}
\end{itemize}

\paragraph{Notes}\label{notes-8}

\begin{itemize}
\tightlist
\item
  Получение кандидатов быстрое.
\end{itemize}

\subsubsection{nearDuplicatesFullScan}\label{nearduplicatesfullscan}

\begin{itemize}
\tightlist
\item
  Operation type: \texttt{read/query}
\item
  Parameters:

  \begin{itemize}
  \tightlist
  \item
    \texttt{docs}: 1500, 3000
  \item
    \texttt{seed}: 42
  \item
    \texttt{threshold}: 0.9
  \item
    \texttt{wordsPerDoc}: 16
  \end{itemize}
\end{itemize}

\paragraph{Notes}\label{notes-9}

\begin{itemize}
\tightlist
\item
  Full scan очень дорогой, но и нужен только для сравнения.
\end{itemize}

\subsubsection{nearDuplicatesLsh}\label{nearduplicateslsh}

\begin{itemize}
\tightlist
\item
  Operation type: \texttt{read/query}
\item
  Parameters:

  \begin{itemize}
  \tightlist
  \item
    \texttt{docs}: 6000, 12000
  \item
    \texttt{seed}: 42
  \item
    \texttt{threshold}: 0.9
  \item
    \texttt{wordsPerDoc}: 16
  \end{itemize}
\end{itemize}

\paragraph{Notes}\label{notes-10}

\begin{itemize}
\tightlist
\item
  LSH-поиск существенно быстрее full scan.
\end{itemize}

\subsection{PerfectHashingBenchmark}\label{perfecthashingbenchmark}

{\def\LTcaptype{none} % do not increment counter
\begin{longtable}[]{@{}
  >{\raggedright\arraybackslash}p{(\linewidth - 10\tabcolsep) * \real{0.1304}}
  >{\raggedleft\arraybackslash}p{(\linewidth - 10\tabcolsep) * \real{0.1739}}
  >{\raggedleft\arraybackslash}p{(\linewidth - 10\tabcolsep) * \real{0.1739}}
  >{\raggedleft\arraybackslash}p{(\linewidth - 10\tabcolsep) * \real{0.1739}}
  >{\raggedleft\arraybackslash}p{(\linewidth - 10\tabcolsep) * \real{0.1739}}
  >{\raggedleft\arraybackslash}p{(\linewidth - 10\tabcolsep) * \real{0.1739}}@{}}
\toprule\noalign{}
\begin{minipage}[b]{\linewidth}\raggedright
operation
\end{minipage} & \begin{minipage}[b]{\linewidth}\raggedleft
kind
\end{minipage} & \begin{minipage}[b]{\linewidth}\raggedleft
duration\_s
\end{minipage} & \begin{minipage}[b]{\linewidth}\raggedleft
p95\_cpu\_\%
\end{minipage} & \begin{minipage}[b]{\linewidth}\raggedleft
peak\_rss\_mb
\end{minipage} & \begin{minipage}[b]{\linewidth}\raggedleft
score\_median
\end{minipage} \\
\midrule\noalign{}
\endhead
\bottomrule\noalign{}
\endlastfoot
buildIndex & build & 326.361 & 0.300 & 3027.984 & 3386.281 \\
containsHit & read/query & 323.100 & 0.200 & 331.641 & 15.627 \\
containsMiss & read/query & 323.184 & 0.300 & 331.859 & 14.979 \\
getHit & read/query & 323.207 & 0.300 & 330.938 & 13.926 \\
getMiss & read/query & 323.330 & 0.300 & 330.297 & 14.738 \\
\end{longtable}
}

\subsubsection{buildIndex}\label{buildindex-1}

\begin{itemize}
\tightlist
\item
  Operation type: \texttt{build}
\item
  Parameters:

  \begin{itemize}
  \tightlist
  \item
    \texttt{n}: 200000, 800000
  \item
    \texttt{seed}: 42
  \end{itemize}
\end{itemize}

\paragraph{Notes}\label{notes-11}

\begin{itemize}
\tightlist
\item
  Индекс строится дорого и с большим разбросом: avgt median
  \textasciitilde38.6 ms/op (в переводе из 38642 us/op).
\item
  Фикс: сократить число повторений при подборе 2-го уровня, а также
  можно заранее сделать pre-size массивы 2-го уровня по размеру
  bucket\^{}2.
\end{itemize}

\subsubsection{containsHit}\label{containshit}

\begin{itemize}
\tightlist
\item
  Operation type: \texttt{read/query}
\item
  Parameters:

  \begin{itemize}
  \tightlist
  \item
    \texttt{n}: 200000, 800000
  \item
    \texttt{seed}: 42
  \end{itemize}
\end{itemize}

\paragraph{Notes}\label{notes-12}

\begin{itemize}
\tightlist
\item
  Просмотр быстрый и стабильный.
\end{itemize}

\subsubsection{containsMiss}\label{containsmiss}

\begin{itemize}
\tightlist
\item
  Operation type: \texttt{read/query}
\item
  Parameters:

  \begin{itemize}
  \tightlist
  \item
    \texttt{n}: 200000, 800000
  \item
    \texttt{seed}: 42
  \end{itemize}
\end{itemize}

\paragraph{Notes}\label{notes-13}

\begin{itemize}
\tightlist
\item
  Miss сопоставим с hit.
\end{itemize}

\subsubsection{getHit}\label{gethit-1}

\begin{itemize}
\tightlist
\item
  Operation type: \texttt{read/query}
\item
  Parameters:

  \begin{itemize}
  \tightlist
  \item
    \texttt{n}: 200000, 800000
  \item
    \texttt{seed}: 42
  \end{itemize}
\end{itemize}

\paragraph{Notes}\label{notes-14}

\begin{itemize}
\tightlist
\item
  getHit чуть хуже containsHit из-за возврата значения.
\item
  Фикс: хранить value в примитивном массиве/структуре.
\end{itemize}

\subsubsection{getMiss}\label{getmiss-1}

\begin{itemize}
\tightlist
\item
  Operation type: \texttt{read/query}
\item
  Parameters:

  \begin{itemize}
  \tightlist
  \item
    \texttt{n}: 200000, 800000
  \item
    \texttt{seed}: 42
  \end{itemize}
\end{itemize}

\paragraph{Notes}\label{notes-15}

\begin{itemize}
\tightlist
\item
  Фикс: унифицировать ветку miss между get/contains для одинакового
  поведения и меньшего кода.
\end{itemize}

\end{document}
